
%%%%%%%%%%%%%%%%%%%%%%%%%%%%%%%%%%%%%%%%%%%%%%%%%%%%%%%%%%%%%%%%%%%%%%%%%%
%
%              582-2
%
%%%%%%%%%%%%%%%%%%%%%%%%%%%%%%%%%%%%%%%%%%%%%%%%%%%%%%%%%%%%%%%%%%%%%%%%%%

\section{Completely Simple Semigroups}
\label{2.2}

     Let us again consider a semigroup $S = X \times G \times Y$ where for
$(x_1,g_1,y_1),(x_2,g_2,y_2) \in S$
\[
     (x_1,g_1,y_1)(x_2,g_2,y_2) = (x_1,g_1 \phi(y_1x_2)g_2, y_2)
\]
 where $\phi: Y \times X \to G$ is arbitrary.  (We will usually write
$\phi(y,x) = (yx).$  Notice the following:
\begin{enumerate}
\item[(1)] Each set of the form $\{ x_1\} \times G \times \{ y_1\}$ is
isomorphic to $G$ under the isomorphism
\[
     \aligned
     g &\mapsto (x_1,g(y_1x_1)^{-1}, y_1)               \\
     \hbox{or } \quad g &\mapsto
          (x_1,(y_1x_1)^{-1}g, y_1).
     \endaligned
\]
(Show this.)

\item[(2)] Note that the local idempotent is of the form
$(x,(yx)^{-1},y)$.  In fact, suppose $(x,h,y)(x,h,y) = (x,h,y)$
\[
     \alignedat{2}
     \hbox{then}& &\quad &h(yx)h = h                    \\
-     \hbox{or}& &\quad &h(yx) = e                       \\
     & &\quad &h = (yx)^{-1}.
     \endalignedat
\]

\item[(3)] One can picture $X \times G \times Y$ as a union of isomorphic
copies of $G$ with a different ``center'' in each box (``center'' is
the local identity).
\[
\setlength{\tabcolsep}{12pt}
\renewcommand{\arraystretch}{1.5}
\begin{tabular}{c|c|c|c|c|c|} \hline
	&$y_1$  &$y_2$  &$y_3$  &$y_4$  &$y_5$ \\ \hline
$x_1$   & & & & & \\ \hline
$x_2$   & & & & & \\ \hline
$x_3$   & & & & & \\ \hline
$x_4$   & & & & & \\ \hline
\end{tabular}
\]


\runinhead{Homework.}\runinhead{Homework} 

\item[(4)] Sets of the form $\{ x\} \times G \times Y$ form minimal right ideals, while minimal left ideals are $X \times G \times \{ y\}$.

\runinhead{Homework.} 
Show that if $L$ is a left ideal of $S$.  
Then $L = X \times G \times Y$, where $Y_1 \subset Y$.  
(A similar result holds for right ideals.)

\item[(5)] Since the local groups of the form $\{ x\} \times G \times \{ Y\}$
are disjoint from each other (that is, $\{ x_1\} \times G \times \{ y_1\} \cap
\{ x_2\} \times G \times \{ y_2\} = \emptyset$ iff $x_1 \not= x_2$ or
$y_1 \not= y_2$), then these form a partition of $S$.
\[
     S = \underset{y_j \in Y}{\bigcup_{x_i \in X}}
          \{ x_i\} \times G \times \{ y_j\} \equiv \bigcup_{i,j} G_{ij}
\]
Also note that $G_{ij} \cdot G_{mn} = G_{in}$.  
That is, the partition $\{ G_{ij}\}$ is preserved under the semigroup operation.  
If we call $\rho$ the equivalence relation that generates $\{ G_{ij}\}$ then

\runinhead{Homework.} 
$S/\rho$ is a quotient semigroup isomorphic to $x \times \{ e\} \times Y$.

\end{enumerate}

\begin{proposition}  	\label{prop:2.1}
Let $S = X \times G \times Y$.  
Then $S$ is simple.
\end{proposition}

\begin{proof}  Suppose $I \subsetneq S$ was a nontrivial ideal of
$S$.  Then for all $a \in I$
\[
     \gathered
     Sa \subset I                                       \\
     \Rightarrow Sa S \subset IS \subset I
     \endgathered
\]
Let $a = (x,g,y)$ be arbitrary.  Then
\[
     \gathered
     (X \times G \times Y)(x,g,y)(X \times G \times Y)                  \\
     = X \times G(Yx)g(yX)G \times Y
     \endgathered
\]
but $GhG = G$ for {\it any} $h \in G$.  Thus $SaS = X \times G \times Y$ and
$I = S$ which implies that $S$ is simple.
\qed
\end{proof}

\begin{definition} 	\label{defn:2.4}
An idempotent $e$ is called primitive iff $e$ is minimal w.r.t. the partial order $e \le f \Leftrightarrow ef = fe = e$.
\end{definition}

\begin{definition} 	\label{defn:2.5}
A semigroup $S$ is called completely simple iff $S$ is simple and contains a primitive idempotent.
\end{definition}

\begin{example}		\label{exm:2.5}  
\[
     S = \left\{ \begin{bmatrix}  a &0 \\ b &1 \end{bmatrix}
          \,\bigg\|\, a,b > 0 \right\}                 
\]
Observe that if 
\[
        \begin{bmatrix}  a_1 &0 \\ b_1 &1 \end{bmatrix} \
        \begin{bmatrix} a_1 &0  \\ b_1 &1 \end{bmatrix} =
        \begin{bmatrix} a_1 &0  \\ b_1 &1 \end{bmatrix},    
\]
then $a^2_1 = a_1$ and $a_1b_1 + b_1 = b_1$.
But $a^2_1=a_1$ implies that $a_1=1$.
Thus $2b_1=b_1$, which forces $b_1 = 0$.
Thus $S$ contains no idempotent.
\end{example}

\runinhead{Homework.} 
Show that the semigroup $S$ in Example \ref{exm:2.5} is simple.

\runinhead{Homework.} 
Show that every finite semigroup contains an idempotent.
